% GENERATED FILE. Edit canonical lessons, then run npm run generate:books.
% canonical-source-hash: fnv1a64:46ef2edbc010ef84
% canonical-lessons: ML-C538-sukham-prapikkuka, ML-C538-paricarikkuka, ML-C538-vilceyar, ML-C538-vedanikkuka, ML-C538-bodham-ketuka

\chapter{Injury and Recovery}
\label{ch:ml-injury-and-recovery}
\hlchaptermodality{538}

\begin{chapteropening}
\textbf{By the end of this chapter:} \emph{I can say to get well, to tend, a wheelchair, to hurt and to pass out.}

Complete the last lesson of chapter 538: I can say to get well, to tend, a wheelchair, to hurt and to pass out.
\end{chapteropening}

\section[\texorpdfstring{\ml{സുഖം} \ml{പ്രാപിക്കുക} (sukhaṁ prāpikkuka)}{സുഖം പ്രാപിക്കുക (sukhaṁ prāpikkuka)}]{\ml{സുഖം} \ml{പ്രാപിക്കുക} (sukhaṁ prāpikkuka) — to get well; to recover}

\label{lesson:ML-C538-sukham-prapikkuka}



\begin{quote}
Before the new one: say the Malayalam for to snore, then the Malayalam for to perspire.
\end{quote}

\subsection*{You'll want to know: \ml{സുഖം} \ml{പ്രാപിക്കുക}}
\textbf{\ml{സുഖം} \ml{പ്രാപിക്കുക}} — \emph{sukhaṁ prāpikkuka} — \textquotedblleft{}to get well; to recover\textquotedblright{}.

\begin{quote}
\ml{ഓപ്പറേഷന്} \ml{ശേഷം} \ml{അച്ഛൻ} \ml{വേഗം} \ml{സുഖം} \ml{പ്രാപിച്ചു}.
\end{quote}

\begin{quote}
\emph{ōppaṟēṣanŭ śēṣaṁ acchan vēgaṁ sukhaṁ prāpiccu.}
\end{quote}

\begin{quote}
Father recovered quickly after the operation.
\end{quote}

Past: \ml{സുഖം} \ml{പ്രാപിച്ചു}, \emph{sukhaṁ prāpiccu}.

Formal, common in hospital talk and the news; in speech one says \ml{അസുഖം} \ml{മാറി}.

\begin{cousinweb}
\textbf{\ml{സുഖം} \ml{പ്രാപിക്കുക}} (\emph{sukhaṁ prāpikkuka}) is built from \ml{സുഖം}, from Sanskrit \emph{sukha}, \textquotedblleft{}pleasant, easy\textquotedblright{} (CDIAL 13451), and \ml{പ്രാപിക്കുക}, built on the Sanskrit verb \emph{prāp}, \textquotedblleft{}to reach, to obtain\textquotedblright{}, whose noun \emph{prāpti} means acquisition (CDIAL 8946): to reach well-being.
\end{cousinweb}

\begin{grammarlens}[title={What you've built}]
One more word for describing what you lived and felt.
\end{grammarlens}

\subsection*{Your turn}
\begin{itemize}
\raggedright
  \item \emph{Say it:} \emph{sukhaṁ prāpikkuka}
  \item \emph{Say it:} \emph{sukhaṁ prāpikkuka}, once more
  \item \emph{Say it:} \emph{ōppaṟēṣanŭ śēṣaṁ acchan vēgaṁ sukhaṁ prāpiccu.} — Father recovered quickly after the operation.
  \item \emph{Recall:} say the Malayalam for to snore, then the Malayalam for to perspire, then say \emph{sukhaṁ prāpikkuka} again
\end{itemize}

\begin{tcolorbox}[breakable,colback=teal!4,colframe=teal!35!black,title={Before you move on}]
What does \ml{സുഖം} \ml{പ്രാപിക്കുക} mean? (\textquotedblleft{}To get well; to recover\textquotedblright{}.) Say it once more.
\end{tcolorbox}
\section[\texorpdfstring{\ml{പരിചരിക്കുക} (paricarikkuka)}{പരിചരിക്കുക (paricarikkuka)}]{\ml{പരിചരിക്കുക} (paricarikkuka) — to tend (the sick); to nurse}

\label{lesson:ML-C538-paricarikkuka}



\begin{quote}
Before the new one: say the Malayalam for to perspire, then the Malayalam for to get well.
\end{quote}

\subsection*{You'll want to know: \ml{പരിചരിക്കുക}}
\textbf{\ml{പരിചരിക്കുക}} — \emph{paricarikkuka} — \textquotedblleft{}to tend (the sick); to nurse\textquotedblright{}.

\begin{quote}
\ml{ആശുപത്രിയിൽ} \ml{ചേച്ചി} \ml{അമ്മയെ} \ml{പരിചരിച്ചു}.
\end{quote}

\begin{quote}
\emph{āśupatriyil cēcci ammaye paricariccu.}
\end{quote}

\begin{quote}
At the hospital my elder sister looked after Mother.
\end{quote}

Past: \ml{പരിചരിച്ചു}, \emph{paricariccu}.

The noun is \ml{പരിചരണം}, care. The one cared for takes -\ml{എ}.

\begin{cousinweb}
\textbf{\ml{പരിചരിക്കുക}} (\emph{paricarikkuka}) is built from Sanskrit \emph{paricaraṇa}, \textquotedblleft{}going about; an attendant\textquotedblright{} (Monier-Williams), with the Malayalam verb ending -\ml{ിക്കുക}.
\end{cousinweb}

\begin{grammarlens}[title={What you've built}]
One more word for describing what you lived and felt.
\end{grammarlens}

\subsection*{Your turn}
\begin{itemize}
\raggedright
  \item \emph{Say it:} \emph{paricarikkuka}
  \item \emph{Say it:} \emph{paricarikkuka}, once more
  \item \emph{Say it:} \emph{āśupatriyil cēcci ammaye paricariccu.} — At the hospital my elder sister looked after Mother.
  \item \emph{Recall:} say the Malayalam for to perspire, then the Malayalam for to get well, then say \emph{paricarikkuka} again
\end{itemize}

\begin{tcolorbox}[breakable,colback=teal!4,colframe=teal!35!black,title={Before you move on}]
What does \ml{പരിചരിക്കുക} mean? (\textquotedblleft{}To tend (the sick); to nurse\textquotedblright{}.) Say it once more.
\end{tcolorbox}
\section[\texorpdfstring{\ml{വീൽചെയർ} (vīlceyar)}{വീൽചെയർ (vīlceyar)}]{\ml{വീൽചെയർ} (vīlceyar) — a wheelchair}

\label{lesson:ML-C538-vilceyar}



\begin{quote}
Before the new one: say the Malayalam for to get well, then the Malayalam for to tend.
\end{quote}

\subsection*{You'll want to know: \ml{വീൽചെയർ}}
\textbf{\ml{വീൽചെയർ}} — \emph{vīlceyar} — \textquotedblleft{}a wheelchair\textquotedblright{}.

\begin{quote}
\ml{ശസ്ത്രക്രിയയ്ക്ക്} \ml{ശേഷം} \ml{അമ്മ} \ml{കുറച്ച്} \ml{ദിവസം} \ml{വീൽചെയറിൽ} \ml{ആയിരുന്നു}.
\end{quote}

\begin{quote}
\emph{śastrakriyaykkŭ śēṣaṁ amma kuṟaccŭ divasaṁ vīlceyaṟil āyirunnu.}
\end{quote}

\begin{quote}
After the surgery Mother was in a wheelchair for a few days.
\end{quote}

Before an ending it adds -\ml{ഇൽ} straight on: \ml{വീൽചെയറിൽ}.

\begin{cousinweb}
\textbf{\ml{വീൽചെയർ}} (\emph{vīlceyar}) was borrowed from English \emph{wheelchair}.
\end{cousinweb}

\begin{grammarlens}[title={What you've built}]
One more word for describing what you lived and felt.
\end{grammarlens}

\subsection*{Your turn}
\begin{itemize}
\raggedright
  \item \emph{Say it:} \emph{vīlceyar}
  \item \emph{Say it:} \emph{vīlceyar}, once more
  \item \emph{Say it:} \emph{śastrakriyaykkŭ śēṣaṁ amma kuṟaccŭ divasaṁ vīlceyaṟil āyirunnu.} — After the surgery Mother was in a wheelchair for a few days.
  \item \emph{Recall:} say the Malayalam for to get well, then the Malayalam for to tend, then say \emph{vīlceyar} again
\end{itemize}

\begin{tcolorbox}[breakable,colback=teal!4,colframe=teal!35!black,title={Before you move on}]
What does \ml{വീൽചെയർ} mean? (\textquotedblleft{}A wheelchair\textquotedblright{}.) Say it once more.
\end{tcolorbox}
\section[\texorpdfstring{\ml{വേദനിക്കുക} (vēdanikkuka)}{വേദനിക്കുക (vēdanikkuka)}]{\ml{വേദനിക്കുക} (vēdanikkuka) — to hurt; to be painful}

\label{lesson:ML-C538-vedanikkuka}



\begin{quote}
Before the new one: say the Malayalam for to tend, then the Malayalam for a wheelchair.
\end{quote}

\subsection*{You'll want to know: \ml{വേദനിക്കുക}}
\textbf{\ml{വേദനിക്കുക}} — \emph{vēdanikkuka} — \textquotedblleft{}to hurt; to be painful\textquotedblright{}.

\begin{quote}
\ml{ഇന്നലെ} \ml{കളിച്ചതുകൊണ്ട്} \ml{കാൽ} \ml{വേദനിക്കുന്നു}.
\end{quote}

\begin{quote}
\emph{innale kaḷiccatukoṇṭŭ kāl vēdanikkunnu.}
\end{quote}

\begin{quote}
My leg hurts because I played yesterday.
\end{quote}

Past: \ml{വേദനിച്ചു}, \emph{vēdaniccu}.

The part that hurts is the subject: \ml{തല} \ml{വേദനിക്കുന്നു}, my head hurts.

\begin{cousinweb}
\textbf{\ml{വേദനിക്കുക}} (\emph{vēdanikkuka}) is built from \ml{വേദന}, \textquotedblleft{}pain\textquotedblright{}, from Sanskrit \emph{vedanā}, \textquotedblleft{}pain, agony\textquotedblright{} (Monier-Williams), with the Malayalam verb ending -\ml{ിക്കുക}.
\end{cousinweb}

\begin{grammarlens}[title={What you've built}]
One more word for describing what you lived and felt.
\end{grammarlens}

\subsection*{Your turn}
\begin{itemize}
\raggedright
  \item \emph{Say it:} \emph{vēdanikkuka}
  \item \emph{Say it:} \emph{vēdanikkuka}, once more
  \item \emph{Say it:} \emph{innale kaḷiccatukoṇṭŭ kāl vēdanikkunnu.} — My leg hurts because I played yesterday.
  \item \emph{Recall:} say the Malayalam for to tend, then the Malayalam for a wheelchair, then say \emph{vēdanikkuka} again
\end{itemize}

\begin{tcolorbox}[breakable,colback=teal!4,colframe=teal!35!black,title={Before you move on}]
What does \ml{വേദനിക്കുക} mean? (\textquotedblleft{}To hurt; to be painful\textquotedblright{}.) Say it once more.
\end{tcolorbox}
\section[\texorpdfstring{\ml{ബോധം} \ml{കെടുക} (bōdhaṁ keṭuka)}{ബോധം കെടുക (bōdhaṁ keṭuka)}]{\ml{ബോധം} \ml{കെടുക} (bōdhaṁ keṭuka) — to pass out; to faint}

\label{lesson:ML-C538-bodham-ketuka}



\begin{quote}
Before the new one: say the Malayalam for a wheelchair, then the Malayalam for to hurt.
\end{quote}

\subsection*{You'll want to know: \ml{ബോധം} \ml{കെടുക}}
\textbf{\ml{ബോധം} \ml{കെടുക}} — \emph{bōdhaṁ keṭuka} — \textquotedblleft{}to pass out; to faint\textquotedblright{}.

\begin{quote}
\ml{വെയിലത്ത്} \ml{നിന്ന} \ml{ഒരു} \ml{കുട്ടി} \ml{ബോധം} \ml{കെട്ടു} \ml{വീണു}.
\end{quote}

\begin{quote}
\emph{veyilattŭ ninna oru kuṭṭi bōdhaṁ keṭṭu vīṇu.}
\end{quote}

\begin{quote}
A child standing in the sun fainted and fell.
\end{quote}

Past: \ml{ബോധം} \ml{കെട്ടു}, \emph{bōdhaṁ keṭṭu}.

\ml{ബോധം} awareness + \ml{കെടുക} to go out, as a lamp does.

\begin{cousinweb}
\textbf{\ml{ബോധം} \ml{കെടുക}} (\emph{bōdhaṁ keṭuka}) is built from Sanskrit \emph{bodha}, \textquotedblleft{}waking, consciousness\textquotedblright{} (Monier-Williams), and the native verb \ml{കെടുക}, \textquotedblleft{}to be put out, to be spoilt\textquotedblright{} (DEDR 1942).
\end{cousinweb}

\begin{grammarlens}[title={What you've built}]
One more word for describing what you lived and felt.
\end{grammarlens}

\subsection*{Your turn}
\begin{itemize}
\raggedright
  \item \emph{Say it:} \emph{bōdhaṁ keṭuka}
  \item \emph{Say it:} \emph{bōdhaṁ keṭuka}, once more
  \item \emph{Say it:} \emph{veyilattŭ ninna oru kuṭṭi bōdhaṁ keṭṭu vīṇu.} — A child standing in the sun fainted and fell.
  \item \emph{Recall:} say the Malayalam for a wheelchair, then the Malayalam for to hurt, then say \emph{bōdhaṁ keṭuka} again
\end{itemize}

\begin{tcolorbox}[breakable,colback=teal!4,colframe=teal!35!black,title={Before you move on}]
What does \ml{ബോധം} \ml{കെടുക} mean? (\textquotedblleft{}To pass out; to faint\textquotedblright{}.) Say it once more.
\end{tcolorbox}
